\section{Problem Formulation}
\label{sec:problem}

This section formalizes the carbon-aware FaaS scheduling problem.
\S\ref{sec:problem:statement} sets up the entities, decision variables,
and constraints. \S\ref{sec:problem:online} frames the optimization as
an online problem and discusses why it admits neither a tractable
offline optimum nor a competitive-ratio bound, motivating the
empirical evaluation strategy in \S\ref{sec:evaluation}.

\subsection{Problem Statement}
\label{sec:problem:statement}

\paragraph{Regions.} A set $\mathcal{R} = \{r_1, \ldots, r_R\}$ of
geographic regions. Each region $r$ has capacity $K_r$, a time-varying
grid carbon intensity $C_r: \mathbb{R}_{\ge 0} \to \mathbb{R}_{>0}$
in gCO$_2$eq/kWh, a Power Usage Effectiveness $\mathrm{PUE}_r \ge 1$,
and network round-trip latencies $\rho(r, r') \ge 0$ to every other
region.

\paragraph{Functions.} A set $\mathcal{F}$ of deployed functions, each
$f \in \mathcal{F}$ characterized by expected execution time
$\tau_e(f)$, cold-start duration $\tau_c(f)$, active power $P_a(f)$,
idle power $P_i(f)$, latency class
$\ell(f) \in \{\Interactive, \Deferrable, \Background\}$, and SLA
deadline $D(f)$. The latency class determines the scheduler's degree
of freedom: \Interactive{} invocations must execute immediately in
their home region; \Deferrable{} invocations may be deferred up to a
class-dependent horizon and routed within a tight RTT budget;
\Background{} invocations may use their full deadline and any region.

\paragraph{Invocations.} Function invocations arrive as a stream
$\mathcal{I} = \langle i_1, i_2, \ldots \rangle$ ordered by arrival
time. Each $i$ specifies a function $f(i)$, an arrival time $a(i)$, a
home region $h(i)$, and a realised runtime $\tau(i)$ drawn from
$f(i)$'s execution-time distribution. Crucially, $\tau(i)$ is
\emph{not} known to the scheduler at decision time; only the
expected runtime $\tau_e(f(i))$ is. The same applies to future
arrivals: at time $a(i)$ the scheduler sees $i$ but knows nothing
about $i_k$ for $k$ with $a(i_k) > a(i)$ beyond what an arrival-rate
estimator can infer.

\paragraph{Warm pools.} For each $(r, f) \in \mathcal{R} \times
\mathcal{F}$, the platform maintains a non-negative integer
warm-container count $W_{r,f}(t)$ representing containers that have
previously executed $f$ in $r$ and remain alive awaiting reuse. A
warm container is evicted after a TTL of $T_w$ seconds of inactivity.

\paragraph{Decisions.} For each invocation $i$, the scheduler produces
$\pi(i) = (r(i), s(i))$: the execution region $r(i)$ and the
scheduled start time $s(i)$. The secondary outcome
$w(i) = \mathbf{1}[W_{r(i), f(i)}(s(i)) > 0]$ --- whether $i$ finds a
warm container at $r(i)$ at time $s(i)$ --- is observed rather than
chosen; if $w(i) = 0$, the invocation pays $\tau_c(f(i))$ in
start-up latency and energy.

\paragraph{Per-invocation cost.} The carbon cost of $i$ under
$\pi(i) = (r, s)$ is
\begin{equation*}
\mathrm{Carbon}(i) = \tfrac{\mathrm{PUE}_r}{3.6 \times 10^6}\,C_r(s)
\cdot P_a(f(i))\big[\tau(i) + (1-w(i))\tau_c(f(i))\big].
\end{equation*}
End-to-end latency is
\begin{equation*}
\mathrm{Lat}(i) = (s - a(i)) + \rho(h(i), r) + (1-w(i))\tau_c(f(i)) + \tau(i).
\end{equation*}

\paragraph{System cost.}
\label{sec:problem:cost}
Aggregate carbon over horizon $[0, T]$ sums
per-invocation execution carbon and warm-pool idle carbon. The exact
idle carbon integrates the time-varying intensity over each idle
interval:
\begin{align}
\mathrm{TotalCarbon}([0,T]) ={}& \sum_{i: a(i) \in [0,T]} \mathrm{Carbon}(i) \nonumber \\
& + \sum_{r,f} \tfrac{\mathrm{PUE}_r}{3.6 \times 10^6} P_i(f) \int_0^T W_{r,f}(t)\,C_r(t)\,dt,
\label{eq:totalcarbon}
\end{align}
where $W_{r,f}(t)$ is the number of warm-but-idle containers for
function $f$ in region $r$ at time $t$, and $C_r(t)$ is the
instantaneous carbon intensity. For analytical tractability in
Lemma~\ref{lemma:tradeoff} and in the \textsc{ScoreSlots} heuristic,
we approximate $C_r(t)$ over an idle interval by its mean
$\bar{C}_{r}$. This approximation is exact when warm-pool idle time
is distributed uniformly over the horizon (the typical case for
steady FaaS load, where the warm pool is continuously cycled);
deviation arises only if idle time correlates with carbon intensity.
Because GreenFaaS's spatial routing concentrates load in
low-carbon regions during high-carbon periods, any residual
correlation is conservative (it would slightly \emph{over}-estimate
idle carbon in the regions GreenFaaS avoids). The simulator
(\S\ref{sec:simulator}) reports results under this mean-intensity
approximation; the scheduler \emph{ranking} is what the paper claims,
and the power-model robustness check (\S\ref{sec:discussion}) confirms
the ranking is insensitive to the idle-energy magnitude.

\paragraph{Constraints.} A schedule $\pi$ must respect:
(1) \emph{Capacity:} at every $t$, the count of in-flight invocations
in region $r$ is at most $K_r$. (2) \emph{SLA-feasibility in
expectation:} $(s(i) - a(i)) + \rho(h(i), r(i)) + (1 - \hat w(i))
\tau_c(f(i)) + \tau_e(f(i)) \le D(f(i))$, where $\hat w(i)$ is the
scheduler's warm-pool prediction. (3) \emph{Class-dependent shifting
limits:} \Interactive{} $\Rightarrow s(i) = a(i)$ and $r(i) = h(i)$;
\Deferrable{} $\Rightarrow s(i) \le a(i) + \Delta^{\mathrm{D}}$ and
$\rho(h(i), r(i)) \le \rho^{\mathrm{D}}$; \Background{} $\Rightarrow$
analogous with class-specific bounds $\Delta^{\mathrm{B}}$,
$\rho^{\mathrm{B}}$. (4) \emph{Causality:} $s(i) \ge a(i)$.

\paragraph{Objective.} Minimise total carbon over $[0, T]$:
\begin{equation}
\pi^* = \arg\min_{\pi}\;\mathrm{TotalCarbon}([0,T];\,\pi)
\quad \text{s.t. (1)--(4)}.
\label{eq:objective}
\end{equation}

\subsection{Online Optimisation Framing}
\label{sec:problem:online}

Two structural features distinguish this problem from prior
carbon-aware scheduling work.

\paragraph{Online, non-clairvoyant.} Decisions $\pi(i)$ must be
committed at time $a(i)$, before any subsequent invocation is
observed. The scheduler has access to past arrivals, the current state
$\{W_{r,f}(t), K_r - \mathrm{in\text{-}flight}_r(t)\}$, and a forecast
$\hat{C}_r$ of future carbon intensity over a finite horizon
(empirically up to 24h with degrading accuracy;
see~\S\ref{sec:evaluation:forecast}). It does \emph{not} see future
arrivals. We further restrict to \emph{non-clairvoyant} schedulers:
the realised runtime $\tau(i)$ is not revealed until execution
completes; the scheduler uses $\tau_e(f(i))$ as its working estimate.
Prior carbon-aware schedulers for batch workloads sometimes assume
per-job runtime is known up front \citep{hanafy2023carbonscaler}; we
cannot.

\paragraph{Non-additive cost structure.} The system cost
(Eq.~\ref{eq:totalcarbon}) is \emph{not} separable into
per-invocation contributions because of the warm-pool idle-energy
term. A scheduler's decision $\pi(i)$ affects not only
$\mathrm{Carbon}(i)$ but also $W_{r,f}(t)$ for all $t \ge s(i)$,
which in turn governs whether subsequent invocations of $f$ in $r$
pay a cold start and how much idle energy accrues. Decisions are
thus \emph{coupled across invocations} through the warm-pool state.
This is the structural reason why batch-oriented online algorithms
--- which typically assume separable per-job costs --- do not
directly apply.

\paragraph{Why no competitive-ratio bound.} Even with full knowledge
of future arrivals, the offline problem contains
generalized-assignment-like structure: invocations must be mapped to
(region, time-slot) bins under per-region capacity constraints $K_r$
and per-function deadline constraints $D(f)$, with a bin-dependent
carbon cost. The classical generalized assignment problem is
$\mathbf{NP}$-hard, and the non-separable idle-energy term makes our
cost function supermodular rather than linear in the assignment, so
we expect the offline optimum to be at least as hard; we do not give
a formal reduction, since the offline problem is not the focus of our
contribution. We do not pursue an approximation-ratio analysis for
two reasons. \emph{First}, the cost function combines additive
execution carbon with a non-separable idle-energy integral, which
puts the problem outside the standard frameworks for online
competitive analysis (makespan, machine scheduling, $k$-server).
\emph{Second}, an arbitrarily small adversarial perturbation of the
carbon-intensity functions $C_r$ can make any online scheduler
arbitrarily worse than the offline optimum, so any finite competitive
ratio would have to assume a specific stochastic model of $C_r$ ---
and the empirical literature is clear that real grid carbon is
neither stationary nor i.i.d.\ across regions. We therefore adopt an
\emph{empirical} evaluation strategy (\S\ref{sec:evaluation}); the
analytical contribution of the paper is the \emph{local} trade-off
characterization of \S\ref{sec:tradeoff}, which gives the scheduler a
principled decision rule without claiming global optimality.

\paragraph{Decomposition.} \GreenFaaS{} decomposes the per-invocation
problem into two sub-problems tractable individually. \emph{Region
gating} (long-run, structural): given the function's arrival rate
$\hat\lambda(f, h)$ and current carbon intensities, which subset of
regions could ever be a carbon-cheaper warm-pool host than the home
region? Lemma~\ref{lemma:tradeoff} (\S\ref{sec:tradeoff}) answers
this in closed form. \emph{Per-slot scoring} (local, this-invocation):
given the admitted regions and a deferral horizon, which
(region, start-time) pair minimises expected carbon for this single
invocation, charged appropriately for idle-energy effects
(\S\ref{sec:tradeoff:idle})? \S\ref{sec:algorithm} makes the
algorithm and its complexity precise.
